\hwhat{We measure how far one operator message moves the agents who see it, and whether that number is trackable by day. We test 16 non-holdout goal periods, after a synthetic validation that fixed three estimator biases. \textbf{Round 1b} uses the corrected activity table, because round 1 lost half of all events. It adds the leading-@ nudge target, receiving-call timing and past-only controls with day effects. \emph{Activity:} a nudge adds about one active minute for its target (\#51 0.98 [0.63, 1.37], \#38 0.95 [0.38, 1.54]; round 1: 0.59, $-0.08$). First nudges give 1.21, repeats 0.44, room-mates $\approx0$. The round-1 lull warning reverses (lull 1.74 vs 0.72). \emph{Content:} human messages turn the next statements toward their new content, $+0.024$ (room-mates) and $+0.041$ (named), in both embedding models. \emph{Daily gauge:} \#51 varies by day ($p=0.026$, $R_1=0.48$) without persistence. NE44: nudges buy more under the 5-min pause default (reversed). The holdout is not run.}

\hmodel{Linear response to impulsive fields (01, 11). Each operator message $k$ is a field pulse on its recipients $i$. Spins: activity $n_i(t)\in\{0,1\}$; content $s_i$, a whitened statement embedding. $u^\perp_k$: the message direction projected off the recipient's last 8 statements. $\chi_{\rm act}$: the 30-min integrated activity response (local projection; strata $\alpha_s$, day effects $\delta_d$, kick counts $K$). $\chi_{\rm con}$: post-message alignment minus same-kind placebos $p$. $\chi_{\rm op}(d)$: the day mean of per-kick responses $r_k$.}

\hmath{\vspace{-5pt}\begin{gather*}
Y_i(m)\equiv\textstyle\sum_{\tau=1}^{30}n_i(m+\tau)=\alpha_{s(i,m)}+\delta_d+\sum_c\chi^c_{\rm act}K^c_i(m)+\beta\cdot K^{\rm adj}_i(m)+\varepsilon\\
\chi_{{\rm con},k}=\bar s_i^{\,\rm post}\!\cdot u^\perp_k-\big\langle \bar s_i^{\,\rm post}\!\cdot u^\perp_p\big\rangle_p,\qquad u^\perp=\frac{(1-P_{\rm pre})u}{\|(1-P_{\rm pre})u\|}\\
\chi_{\rm op}(d)=\frac{\sum_{k\in d}w_kr_k}{\sum_{k\in d}w_k},\qquad R_1=1-\frac{\mathrm{Var}_{\rm perm}[\chi_{\rm op}(d)]}{\mathrm{Var}[\chi_{\rm op}(d)]}
\end{gather*}\vspace{-7pt}}

\hobs{../figures/summary_obs.pdf}{(a) \#51 daily nudge gauge: extra active minutes per nudge for the target (day FE, 95\% CIs), with the period mean and band. Day-to-day scatter matches the daily CIs: no trend, no persistence. (b) Period-level $\chi_{\rm op}$. Nudge$\to$activity in the three largest nudge periods (pre-registered model); first vs repeat nudges in \#51 (day FE); human$\to$content in human-dense periods (cosine $\times20$).}

\htelos{Do not watch a daily ``steerability'' number: at village message rates it is a period constant plus noise, and it does not age. One nudge to an idle agent buys about one extra active minute from its receiving call. A first nudge buys 1.2 min and a repeat 0.4 (CI includes 0). Nudging into a swarm-wide lull works even better (1.74 vs 0.72 min), and short scaffold pauses help (NE44). To change \emph{what} an agent writes, name it and add new content ($+0.041$ vs $+0.024$ for room-mates). For analysts: do not condition on future nudges, as H04's isolation rule did.}

\hbottom{Operator messages have a period-level susceptibility (about one active minute per nudge; a small content pull), but daily steerability is not a trackable state.}

\hresults{\footnotesize\setlength{\tabcolsep}{2pt}\begin{tabular}{@{}p{0.25\columnwidth}p{0.57\columnwidth}p{0.15\columnwidth}@{}}
\textbf{Prediction} & \textbf{Round 1 $\to$ round 1b} & \textbf{Verdict}\\
P1 nudge$\to$target & \#51 0.59 $\to$ 0.98 [0.63, 1.37]; \#38 $-0.08\to$ 0.95 [0.38, 1.54]; small periods 5/10 $>0$ & \checkmark{} (count rule $\times$)\\
P2 bystanders $\approx0$ & \#51 0.13 $\to$ $-0.01$; \#38 $-0.54\to-0.09$ & \checkmark\\
P3 human$\to$activity & $\leq0$ / n.s. in regime I (unchanged) & $\times$\\
P4 human$\to$content & 0.026 $\to$ 0.024 (gte 0.024); named 0.046 $\to$ 0.041 & \checkmark\\
P5 nudge$\to$content & \#51 0.003 $\to$ 0.002 [$-0.001$, 0.005] & \checkmark\\
P6 no daily state & \#51 $p$ 0.08 $\to$ 0.026, $R_1$ 0.28 $\to$ 0.48, lag-1 0.08 & $\times$ as written\\
P7 no aging & \#51 $+0.02$ [$-0.03$, 0.07]/day & \checkmark\\
P8 context sawtooth & early 0.66 vs late 1.49 & $\times$\\
P10 family & no Holm-significant lab & \checkmark\\
P11 nulls & \#51 and \#38 now pass & \checkmark\\
Native \#5 dose & pull falls 0.025$\to$0.013 over 1--7 msgs; 8+ noisy & $\times$ (rule)\\
Native NE44 & 12-h 0.57 vs 5-min 0.98 ($\Delta$ $-0.41$) & $\times$ reversed\\
Native NE10 & first nudges ever $-0.42$ (22 kicks) & $\times$ low power\\
\end{tabular}}

\hobsb{../figures/summary_obs_r1b.pdf}{Round 1b. (a) What moved the nudge response in \#51 (blue) and \#38 (orange). Restoring the dropped events carries most of the change; past-only controls bring the level back to $\approx1$ min. (b) Situational $\chi_{\rm op}$ in \#51 (first vs repeat nudges, swarm lull vs not, bystanders) and human content pull ($\times20$).}

\hcaveats{Only \#51 and \#38 resolve the nudge response; 3--5-day periods give unstable CIs. The level depends on the future-kick adjustment: 0.98 past-only vs 1.56 with the A1 regressors of round 1 (day FE). SE-based window reliabilities exceed 0.5 in \#51 but also in \#38, where the calibrated permutation finds no day signal, so we do not count them. Human effects on activity in dense regime-I chat are not identified. Content measures message-specific steering only, and human replies may share topics. The round-1b design (receiving-call timing, past-only, day FE primary) follows the re-evaluation brief, not the pre-registration.}
